
\begin{theorem}[\statusfalse; computer-assisted proof]\label{thm:theta}
At $n=9$ and $t=1/50$, one has $J_9(t)<0$.
\end{theorem}
\begin{proof}
Put $u=\pi m^2e^{4t}$ and define integer polynomials
\[
P_0(u)=2u-3,\qquad P_{k+1}(u)=4uP_k'(u)+(5-4u)P_k(u).
\]
Termwise differentiation gives
\[
\Phi^{(k)}(t)=\sum_{m\ge1}\pi m^2e^{5t-u}P_k(u).
\]
Outward-rounded interval summation, with an explicit geometrically bounded tail, yields
\begin{align*}
\Phi^{(8)}(1/50)&\approx 2.95940368409056694\cdot10^8,\\
\Phi^{(9)}(1/50)&\approx 5.80442287530267756\cdot10^9,\\
\Phi^{(10)}(1/50)&\approx 2.02040519425950258\cdot10^{11},
\end{align*}
and an interval for $J_9(1/50)$ contained strictly in the negative half-line, centered at
\[
-2.61006208371358922\cdot10^{19}.
\]
The source package contains both a high-precision independent evaluator and a Sage \texttt{RealIntervalField} certificate.  The negative margin is many orders of magnitude larger than the certified truncation error.
\end{proof}
